\label{sec:conclusion}

\paragraph{Conclusion.}
This report summarised \texttt{X3D} \cite{feichtenhofer2020x3d}, the family of video networks obtained by expanding a small 2D architecture one axis at a time, and examined how its \texttt{PyTorchVideo} implementation \cite{fan2021pytorchvideo} transfers to isolated sign language recognition on \texttt{AUTSL} \cite{sincan2020autsl}.
Our best configuration, \texttt{X3D\_M} with two frozen blocks, reaches $82.28\%$ top--1, $94.49\%$ top--3 and $96.52\%$ top--5 on the \texttt{AUTSL} test split, at $4.90$ GFLOPs per clip and $3.42$M trainable parameters.
On the first question, full fine--tuning is worth between $39.4$ and $47.5$ top--1 points over a retrained head, but the gain is unevenly distributed: unfreezing one further block recovers most of it, and freezing the first one or two blocks costs nothing.
On the second, the larger variant repays its cost only in part.
The step from \texttt{X3D\_XS} to \texttt{X3D\_S} returns $19.7$ points for $3.2\times$ the GFLOPs, whereas the step to \texttt{X3D\_M} returns $5.8$ for a further $2.4\times$.
Since all three variants share one parameter count, this is the price of sampling a longer clip rather than of a larger model.

\paragraph{Limitations.}
Each configuration was trained once, so differences of one or two points cannot be separated from seed variance.
Only the RGB modality was used, leaving the depth and skeleton streams of \texttt{AUTSL} unexploited, and GFLOPs measures arithmetic rather than latency on an actual edge device.
Another limitation is, the model can only predicts 226 Turkish isolated signs, any thus, fails to predict out–of–distribution data.

\paragraph{Future work.} \texttt{PyTorchVideo}'s \texttt{X3D} is faster than the original implementation \cite{feichtenhofer2020x3d, pyslowfast}, yet still short of what edge deployment demands.
\texttt{PyTorchVideo} also supports \texttt{EfficientX3D} for that purpose, but only \texttt{X3D\_XS} and \texttt{X3D\_S} carry \texttt{Kinetics–400} weights, \href{https://github.com/facebookresearch/pytorchvideo/blob/f3142bb05cdb56af0704ab6f0adfb0c7bbafe4a0/pytorchvideo/models/hub/efficient_x3d_mobile_cpu.py}{see here}.
Since fine--tuning outperforms training from scratch on small datasets \cite{LI2020103853, 7426826}, as applied to sign language recognition in \cite{10627077, töngi2021applicationtransferlearningsign, shuchang2022surveyhumanactionrecognition, han2022efficient}, we leave \texttt{EfficientX3D} to future work for want of pretrained weights.